% ============================================================
% abstract.tex - 摘要与关键词
%
% 硬性要求：摘要页必须独占一整页
% - 页面内只包含"摘要"标题 + 摘要正文 + 关键词
% - 不允许混入目录、正文、页眉等其他内容
% - 摘要正文字数参考比赛规则（国赛 300-500 字，美赛 1 页）
%
% 技术实现：
% - 使用 \thispagestyle{empty} 隐藏摘要页的页眉页脚
% - 用 \vfill 把关键词推到页面底部，视觉上撑满一页
% - 溢出时用 \begin{spacing}{1.25} 或 \small 压缩
% ============================================================

\thispagestyle{empty}

\begin{center}
  {\Large\bfseries\heiti 摘\quad 要}
\end{center}
\vspace{1em}

% ---- 摘要正文 ----
% 结构建议（每段各自独立）：
%   第1段：问题背景 + 整体思路（1-2 句带过）
%   第2段：问题一 —— 方法 + 具体数值结果
%   第3段：问题二 —— 方法 + 具体数值结果
%   第N段：其余问题 ...
%   末段：模型优点 + 局限性 + 推广价值（简短）
%
% 关键：每问必须给出**具体数值结果**，禁止只写"效果良好/结果合理"
%       评委先看摘要，数字就是得分点

本文针对 XXX 问题，建立了 XXX 模型……（示例占位，请替换为实际内容）

针对问题一，采用 XXX 方法，求解得到 XXX = 1256.3（具体数值）……

针对问题二，建立 XXX 模型，结果表明 XXX 比基准方法提升 12.3\%……

（按实际问题数继续段落）

本文模型的优点在于 XXX；不足是 XXX；可推广至 XXX 场景。

\vfill

\noindent\textbf{关键词：} 关键词1\quad 关键词2\quad 关键词3\quad 关键词4
